\documentclass{article}
\usepackage{amsmath,amssymb,amsthm,algorithm,algorithmic}
\usepackage{hyperref}
\newtheorem{theorem}{Theorem}
\newtheorem{lemma}{Lemma}
\newtheorem{assumption}{Assumption}
\newtheorem{definition}{Definition}
\newtheorem{remark}{Remark}
\begin{document}

%%%%%%%%%%%%%%%%%%%%%%%%%%%%%%%%%%%%%%%%%%%%%%%%%%%%%%%%%%%%
\section*{Algorithm Name}
\textbf{Mirror Descent with Bregman Proximal Step} (also called
\emph{Relative Gradient Descent}).  

%%%%%%%%%%%%%%%%%%%%%%%%%%%%%%%%%%%%%%%%%%%%%%%%%%%%%%%%%%%%
\section*{Mathematical Formulation}
Let $f:\mathbb{R}^{n}\rightarrow\mathbb{R}$ be a convex differentiable
objective and let $h:\mathbb{R}^{n}\rightarrow\mathbb{R}$ be a
$\sigma$‑strongly convex prox–function (with respect to a norm
$\|\cdot\|$).  

The Bregman divergence induced by $h$ is
\[
D_{h}(x,y)\;:=\;h(x)-h(y)-\langle\nabla h(y),\,x-y\rangle .
\tag{1}
\]

Given a stepsize sequence $\{\eta_k\}_{k\ge 0}$, the iterates are defined
by the implicit proximal problem
\[
\boxed{
x_{k+1}
\;=\;
\arg\min_{x\in\mathbb{R}^{n}}
\Big\{
\langle \nabla f(x_k),\,x\rangle
+\frac{1}{\eta_k}\,D_{h}(x,x_k)
\Big\}.
}
\tag{2}
\]

When $h(x)=\tfrac12\|x\|_2^{2}$, $D_h$ reduces to the squared Euclidean
distance and (2) becomes the standard (projected) gradient step.

%%%%%%%%%%%%%%%%%%%%%%%%%%%%%%%%%%%%%%%%%%%%%%%%%%%%%%%%%%%%
\section*{Pseudocode}
\begin{algorithm}[H]
\caption{Mirror Descent (Bregman Proximal Step)}
\begin{algorithmic}[1]
\REQUIRE Initial point $x_0\in\mathbb{R}^{n}$, stepsize schedule
$\{\eta_k\}_{k\ge0}>0$, prox‑function $h$ (strongly convex).
\FOR{$k=0,1,2,\dots$}
    \STATE Compute gradient $g_k\gets\nabla f(x_k)$.
    \STATE Solve the proximal subproblem
    \[
    x_{k+1}\gets\arg\min_x\Big\{\langle g_k,x\rangle
    +\frac{1}{\eta_k}D_h(x,x_k)\Big\}.
    \]
\ENDFOR
\ENSURE Final iterate $x_K$ (or averaged iterate).
\end{algorithmic}
\end{algorithm}

The subproblem has a closed‑form solution when $h$ is Legendre:
\[
\nabla h(x_{k+1})\;=\;\nabla h(x_k)-\eta_k\,\nabla f(x_k).
\tag{3}
\]

%%%%%%%%%%%%%%%%%%%%%%%%%%%%%%%%%%%%%%%%%%%%%%%%%%%%%%%%%%%%
\section*{Dry Run Example}
Consider the one‑dimensional quadratic
\[
f(w)=(w-3)^{2},\qquad \nabla f(w)=2(w-3).
\]
Choose the Euclidean prox‑function $h(w)=\tfrac12 w^{2}$, hence
$D_h(u,v)=\tfrac12 (u-v)^{2}$ and $\sigma=1$.
Set $w_{0}=0$ and a constant stepsize $\eta=0.1$.
The update (3) becomes
\[
w_{k+1}=w_{k}-\eta\,\nabla f(w_{k}).
\]

\begin{center}
\begin{tabular}{c|c|c|c}
$k$ & $w_k$ & $\nabla f(w_k)=2(w_k-3)$ & $f(w_k)$\\\hline
0 & $0$ & $-6$ & $9$\\
1 & $0-0.1(-6)=0.6$ & $2(0.6-3)=-4.8$ & $(0.6-3)^{2}=5.76$\\
2 & $0.6-0.1(-4.8)=1.08$ & $2(1.08-3)=-3.84$ & $(1.08-3)^{2}=3.6864$\\
3 & $1.08-0.1(-3.84)=1.464$ & $2(1.464-3)=-3.072$ & $(1.464-3)^{2}=2.368\,$
\end{tabular}
\end{center}
The objective value decreases from $9$ to $2.37$ in three iterations,
illustrating the descent property of (2).

%%%%%%%%%%%%%%%%%%%%%%%%%%%%%%%%%%%%%%%%%%%%%%%%%%%%%%%%%%%%
\section*{Assumptions}
\begin{assumption}[Convexity and Smoothness (Relative Smoothness)]
\label{ass:rel_smooth}
The function $f$ is convex and **$L$‑smooth relative to $h$**, i.e.
\[
f(y)\;\le\;f(x)+\langle\nabla f(x),y-x\rangle
+\;L\,D_{h}(y,x),\qquad\forall x,y\in\mathbb{R}^{n}.
\tag{4}
\]
\end{assumption}
\begin{assumption}[Strong Convexity of the Prox‑function]
\label{ass:h_strong}
The prox‑function $h$ is $\sigma$‑strongly convex w.r.t. the norm
$\|\cdot\|$, i.e.
\[
D_{h}(x,y)\;\ge\;\frac{\sigma}{2}\|x-y\|^{2},
\qquad\forall x,y.
\tag{5}
\]
\end{assumption}
\begin{assumption}[Stepsize]
\label{ass:stepsize}
The stepsizes satisfy $0<\eta_k\le\frac{1}{L}$ for all $k$.
\end{assumption}

The pair $(L,\sigma)$ governs the condition number
$\kappa:=L/\sigma$ of the relative geometry.

%%%%%%%%%%%%%%%%%%%%%%%%%%%%%%%%%%%%%%%%%%%%%%%%%%%%%%%%%%%%
\section*{Main Convergence Theorem}
\begin{theorem}[Convergence Rate of Mirror Descent]
\label{thm:conv}
Let Assumptions~\ref{ass:rel_smooth}--\ref{ass:stepsize} hold and
let $x^{\star}\in\arg\min f$.  For a constant stepsize
$\eta_k\equiv\eta\le 1/L$, the iterates generated by \eqref{2}
satisfy for all $K\ge1$
\[
f(x_{K})-f(x^{\star})
\;\le\;
\frac{D_{h}(x^{\star},x_{0})}{\eta\,K}.
\tag{6}
\]
Consequently the method attains an $O(1/K)$ suboptimality rate.
\end{theorem}

\begin{remark}
If $h$ is Euclidean ($\sigma=1$) the bound reduces to the classical
gradient descent rate $ \frac{\|x_{0}-x^{\star}\|^{2}}{2\eta K} $.
\end{remark}

%%%%%%%%%%%%%%%%%%%%%%%%%%%%%%%%%%%%%%%%%%%%%%%%%%%%%%%%%%%%
\section*{Theoretical Properties}
\begin{enumerate}
\item \textbf{Stability.}  Since $D_h$ is non‑negative and the stepsize is
bounded by $1/L$, the sequence $\{D_h(x^{\star},x_k)\}$ is non‑increasing,
ensuring bounded iterates.
\item \textbf{Hyperparameter Sensitivity.}  The bound \eqref{6} depends
linearly on $1/\eta$.  The optimal choice is $\eta=1/L$, yielding the
tightest $O(1/(\!L K))$ constant.
\item \textbf{Comparison to Baselines.}
\begin{itemize}
\item Standard gradient descent corresponds to $h(w)=\tfrac12\|w\|^{2}$.
\item When $f$ is not smooth in the Euclidean sense but is $L$‑smooth
relative to a non‑quadratic $h$, the present algorithm can exploit the
geometry and achieve the same $O(1/K)$ rate, whereas Euclidean GD may
fail to converge.
\end{itemize}
\end{enumerate}

%%%%%%%%%%%%%%%%%%%%%%%%%%%%%%%%%%%%%%%%%%%%%%%%%%%%%%%%%%%%
\section*{Preliminaries (Definitions and Lemmas)}
\begin{definition}[Three‑Point Identity]
For any $x,y,z$,
\[
\langle\nabla h(z)-\nabla h(y),\,x-y\rangle
= D_{h}(x,y)-D_{h}(x,z)-D_{h}(z,y).
\tag{7}
\]
\end{definition}

\begin{lemma}[Descent Lemma (Relative Smoothness)]
Under Assumption~\ref{ass:rel_smooth},
\[
f(x_{k+1})\le f(x_k)+\langle\nabla f(x_k),x_{k+1}-x_k\rangle
+L\,D_h(x_{k+1},x_k).
\tag{8}
\]
\end{lemma}
\begin{proof}
Directly the statement of relative smoothness with $x=x_k$, $y=x_{k+1}$.
\end{proof}

%%%%%%%%%%%%%%%%%%%%%%%%%%%%%%%%%%%%%%%%%%%%%%%%%%%%%%%%%%%%
\section*{Main Proof (Step‑by‑Step Derivation)}
We prove Theorem~\ref{thm:conv} for a constant stepsize
$\eta_k\equiv\eta\le 1/L$.

\noindent\textbf{Step 1 (Optimality condition of the subproblem).}
The optimality condition of \eqref{2} reads
\[
0\;\in\;\nabla f(x_k)+\frac{1}{\eta}\bigl(\nabla h(x_{k+1})-
\nabla h(x_k)\bigr),
\]
hence
\[
\nabla h(x_{k+1})\;=\;\nabla h(x_k)-\eta\,\nabla f(x_k).
\tag{9}
\]
[Step 1]

\noindent\textbf{Step 2 (Apply three‑point identity).}
Take $x=x^{\star}$, $y=x_k$, $z=x_{k+1}$ in \eqref{7}:
\[
\langle\nabla h(x_{k+1})-\nabla h(x_k),\,x^{\star}-x_k\rangle
= D_h(x^{\star},x_k)-D_h(x^{\star},x_{k+1})-D_h(x_{k+1},x_k).
\tag{10}
\]
[Step 2]

\noindent\textbf{Step 3 (Replace the left‑hand side using (9)).}
From (9),
\[
\langle -\eta\,\nabla f(x_k),\,x^{\star}-x_k\rangle
= D_h(x^{\star},x_k)-D_h(x^{\star},x_{k+1})-D_h(x_{k+1},x_k).
\tag{11}
\]
Dividing by $\eta$ gives
\[
\langle \nabla f(x_k),\,x_k-x^{\star}\rangle
= \frac{1}{\eta}\bigl[D_h(x^{\star},x_k)-D_h(x^{\star},x_{k+1})
-D_h(x_{k+1},x_k)\bigr].
\tag{12}
\]
[Step 3]

\noindent\textbf{Step 4 (Upper bound $f(x_{k+1})-f(x^{\star})$).}
Using convexity of $f$,
\[
f(x_{k+1})-f(x^{\star})\le
\langle \nabla f(x_k),\,x_{k+1}-x^{\star}\rangle .
\tag{13}
\]
Add and subtract $x_k$:
\[
\langle \nabla f(x_k),\,x_{k+1}-x^{\star}\rangle
= \langle \nabla f(x_k),\,x_{k+1}-x_k\rangle
+ \langle \nabla f(x_k),\,x_k-x^{\star}\rangle .
\tag{14}
\]
[Step 4]

\noindent\textbf{Step 5 (Bound the first term via relative smoothness).}
From Lemma (Descent Lemma) \eqref{8} and rearranging,
\[
\langle \nabla f(x_k),\,x_{k+1}-x_k\rangle
\le f(x_{k+1})-f(x_k)+L\,D_h(x_{k+1},x_k).
\tag{15}
\]
[Step 5]

\noindent\textbf{Step 6 (Insert (12) into (14) and use (15)).}
Combining (14), (15) and (12):
\[
\begin{aligned}
f(x_{k+1})-f(x^{\star})
&\le f(x_{k+1})-f(x_k)+L\,D_h(x_{k+1},x_k)
\\
&\quad +\frac{1}{\eta}\bigl[D_h(x^{\star},x_k)-D_h(x^{\star},x_{k+1})
-D_h(x_{k+1},x_k)\bigr].
\end{aligned}
\tag{16}
\]
Cancel $f(x_{k+1})$ on both sides and collect the Bregman terms:
\[
f(x_k)-f(x^{\star})
\ge \Bigl(\frac{1}{\eta}-L\Bigr) D_h(x_{k+1},x_k)
+\frac{1}{\eta}\bigl[D_h(x^{\star},x_{k+1})-D_h(x^{\star},x_k)\bigr].
\tag{17}
\]
[Step 6]

\noindent\textbf{Step 7 (Use the stepsize condition).}
Since $\eta\le 1/L$, the coefficient
$\frac{1}{\eta}-L\ge0$, thus we can drop the non‑negative term
$ (\frac{1}{\eta}-L) D_h(x_{k+1},x_k)$ from the right‑hand side of
\eqref{17}, obtaining the fundamental recursion
\[
f(x_k)-f(x^{\star})
\;\ge\;
\frac{1}{\eta}\bigl[D_h(x^{\star},x_{k+1})-D_h(x^{\star},x_k)\bigr].
\tag{18}
\]
[Step 7]

\noindent\textbf{Step 8 (Telescoping sum).}
Summing \eqref{18} for $k=0,\dots,K-1$ yields
\[
\sum_{k=0}^{K-1}\bigl[f(x_k)-f(x^{\star})\bigr]
\;\ge\;
\frac{1}{\eta}\bigl[D_h(x^{\star},x_{K})-D_h(x^{\star},x_{0})\bigr].
\tag{19}
\]
Since $D_h(\cdot,\cdot)\ge0$, the left‑hand side is bounded above by
$K\bigl[f(x_{K})-f(x^{\star})\bigr]$ (the largest term dominates the sum):
\[
K\bigl[f(x_{K})-f(x^{\star})\bigr]
\;\ge\;
\frac{1}{\eta}\bigl[D_h(x^{\star},x_{0})-D_h(x^{\star},x_{K})\bigr]
\;\ge\;
\frac{1}{\eta}D_h(x^{\star},x_{0}).
\tag{20}
\]
Rearranging gives the desired bound \eqref{6}:
\[
f(x_{K})-f(x^{\star})
\;\le\;
\frac{D_h(x^{\star},x_{0})}{\eta\,K}.
\tag{21}
\]
[Step 8]  This completes the proof. \qed

%%%%%%%%%%%%%%%%%%%%%%%%%%%%%%%%%%%%%%%%%%%%%%%%%%%%%%%%%%%%
\section*{Rate Derivation (With Constants)}
The constant in \eqref{21} is exactly the initial Bregman distance.
If $h$ is $\sigma$‑strongly convex, then by \eqref{5}
\[
D_h(x^{\star},x_{0})\;\le\;\frac{\sigma}{2}\|x_{0}-x^{\star}\|^{2},
\]
so a more explicit Euclidean bound reads
\[
f(x_{K})-f(x^{\star})
\;\le\;
\frac{\sigma}{2\eta\,K}\,\|x_{0}-x^{\star}\|^{2}.
\tag{22}
\]
Choosing the maximal admissible stepsize $\eta=1/L$ yields
\[
f(x_{K})-f(x^{\star})
\;\le\;
\frac{L\sigma}{2}\,\frac{\|x_{0}-x^{\star}\|^{2}}{K}
\;=\;O\!\Bigl(\frac{1}{K}\Bigr).
\]

%%%%%%%%%%%%%%%%%%%%%%%%%%%%%%%%%%%%%%%%%%%%%%%%%%%%%%%%%%%%
\section*{Special Cases}
\begin{enumerate}
\item \textbf{Euclidean case.}
$h(x)=\tfrac12\|x\|^{2}$, $D_h=\tfrac12\|x-y\|^{2}$, $\sigma=1$.
The recursion \eqref{18} becomes the classic gradient descent proof,
and \eqref{6} reduces to
\[
f(x_{K})-f(x^{\star})\le\frac{\|x_{0}-x^{\star}\|^{2}}{2\eta K}.
\]

\item \textbf{Entropy prox (probability simplex).}
Take $h(x)=\sum_{i}x_i\log x_i$ on the simplex.
The Bregman divergence is the Kullback–Leibler divergence.
The algorithm coincides with the \emph{Exponentiated Gradient} method,
and the same $O(1/K)$ guarantee holds under relative smoothness
with respect to the KL divergence.

\item \textbf{Accelerated variant.}
If one employs a Nesterov‑type extrapolation on the dual space,
the rate improves to $O(1/K^{2})$ under the same assumptions
(see e.g., Beck–Teboulle, “Mirror Descent and Dual Averaging”).
\end{enumerate}

%%%%%%%%%%%%%%%%%%%%%%%%%%%%%%%%%%%%%%%%%%%%%%%%%%%%%%%%%%%%
\section*{Discussion}
The proof above shows that the Bregman‑proximal update (2) inherits the
classical $O(1/K)$ convergence of gradient descent **without** requiring
global Lipschitz continuity of $\nabla f$ in the Euclidean norm.  The
relative smoothness condition \eqref{4} is substantially weaker and
covers many important problems (e.g., logistic regression with
entropy distance).  The stepsize bound $\eta\le1/L$ is tight: larger
stepsizes would break the non‑negativity of the term
$(\frac{1}{\eta}-L)D_h$ in \eqref{17} and may cause divergence.

Because the bound depends only on the initial Bregman distance,
choosing a prox‑function that captures the geometry of the feasible
set can dramatically reduce the constant factor, leading to faster
practical convergence even though the asymptotic rate remains $O(1/K)$.

\end{document}